\chapter{The Mechanical Origins of Electromagnetism}


%----------------------------------------------------------------------
\section{Introduction: Fields as Geometric Consequences}
%----------------------------------------------------------------------

In 20th-century physics, the electric and magnetic fields are fundamental components of reality, governed by Maxwell's Equations and quantised in QED.  SDT offers a more fundamental explanation: the electric and magnetic fields are the \textbf{radial and helical components of the propagated pressure field} created by a spinning electron vortex.


%----------------------------------------------------------------------
\section{The Electric Field as a Radial Displacement}
%----------------------------------------------------------------------

\textbf{Mechanism:} The electron vortex, by its existence, displaces the spation medium, creating a static, radial pressure gradient.  This persistent pressure gradient \emph{is} the electric field.

\textbf{Coulomb's Law from SDT:} The force between two charged particles is a direct interaction between their respective pressure fields:
\begin{equation}\label{eq:coulomb-sdt}
  F_{\text{electric}} = \frac{c^2 R_e}{k_e^2} \cdot \frac{m_e}{r^2}
\end{equation}

where $R_e$ and $k_e$ pertain to the source electron.  This reproduces the correct $1/r^2$ dependence and magnitude.


%----------------------------------------------------------------------
\section{The Magnetic Field as a Helical Wake}
%----------------------------------------------------------------------

\textbf{Mechanism:} The electron vortex is spinning.  This spin drags the local spation medium, creating a continuous, swirling, \textbf{helical ``wake''} around its rotation axis.  This persistent helical flow \emph{is} the magnetic field.

\textbf{The Magnetic Dipole:}
\begin{equation}\label{eq:B-dipole}
  |\mathbf{B}(r)| \propto (k_e \cdot c) \cdot \left(\frac{R_e}{r}\right)^3 \sin\theta
\end{equation}

where $\theta$ is the angle from the spin axis.  The $r^{-3}$ dipole decay follows directly from the vortex geometry.


%----------------------------------------------------------------------
\section{The Magnetic Properties of Matter}
%----------------------------------------------------------------------

The difference between magnetic and non-magnetic materials is determined by how their electron vortices align and allocate their movement budgets.

\subsection{Diamagnetism (Paired Vortices)}

Spin-paired electrons have opposite helical wakes that cancel: $v_{\text{magnetic,net}} = 0$.  An external field induces slight counter-circulation as vortices resist the external flow---Lenz's Law.  Weak repulsion.

\subsection{Paramagnetism (Unpaired Vortices)}

Unpaired electrons have net magnetic dipoles.  An external $B$-field forces each vortex to allocate budget to $v_{\text{magnetic}}$ for precession into alignment.  The degree of alignment is a competition between alignment energy ($\propto B$) and thermal budget ($k_B T$).

\subsection{Ferromagnetism (Synchronised, Locked Vortices)}

In certain crystal lattices, neighbouring electron vortices can \textbf{geometrically mesh}, locking their helical wakes into a single, coherent, large-scale flow.  This is a minimum-energy state representing a collective budget allocation.  Below the Curie temperature, this geometric lock is stable, creating a powerful permanent field.

\textbf{The Curie temperature} is the point where the thermal budget ($v_{\text{thermal}}$) becomes large enough to break these geometric locks.


%----------------------------------------------------------------------
\section{Electric Current and Faraday's Law of Induction}
%----------------------------------------------------------------------

\subsection{Electric Current as Coherent Vortex Flow}

An electric current is the linear, drifting motion ($v_{\text{drift}}$) of electron vortices through a conductor.  As these spinning vortices move in a line, their individual helical wakes superimpose to create a large-scale, cylindrical flow of the spation medium around the wire.  This is the magnetic field of a current.

The \textbf{right-hand rule} is the natural geometric result of combining linear motion with intrinsic spin.

\subsection{Faraday's Law as Budget Reallocation}

A \emph{changing} magnetic field is a \emph{changing} helical flow pattern.  This change creates a pressure gradient that exerts a direct ``push'' on free electron vortices within a conductor.

\textbf{Mechanism:} To respond to this changing external field, electron vortices must reallocate their budget.  This forced reallocation manifests as coherent linear motion ($v_{\text{linear}}$)---an induced electric current.

\begin{equation}\label{eq:faraday}
  \boxed{\varepsilon = -\frac{d\Phi_{\text{displacement}}}{dt}}
\end{equation}

This provides a direct, mechanical origin for Faraday's Law.


%----------------------------------------------------------------------
\section{Conclusion}
%----------------------------------------------------------------------

SDT provides a complete, mechanical, and deterministic origin for all electromagnetic phenomena:
\begin{itemize}
  \item The \textbf{electric field}: a radial pressure gradient.
  \item The \textbf{magnetic field}: a dynamic, helical wake.
  \item \textbf{Magnetic materials}: geometric alignment and budget allocation of vortices.
  \item \textbf{Currents and induction}: dynamics of vortex flow and budget reallocation.
\end{itemize}

Electricity and magnetism are not separate forces.  They are the inseparable radial and rotational consequences of a single entity: the spinning electron vortex.
